\section{结论}
\label{sec:conclusion}

本文把严格入口预算下的问题适配写成对一个已有视觉集合的有限更新。BELA 用两个视觉深度上的显著性与空间原型新颖性确定初始身份，再用问题嵌入代理提出交换，并用全体视觉位置的最近代表覆盖变化决定这次交换能否发生。进入 Decoder 后，第一次缩减按问题响应相对浅层记录的秩漂移排序，第二次缩减按问题到视觉的注意力保留。初始化与更新使用不同的集合函数，它们通过被传递的索引集合连接，而不是对同一个目标做两轮求解。

在 LLaVA-1.5-7B、$W{=}1024$ 的固定日程上，这一入口更新提高了六任务归一化平均，并在 GQA 与 POPE 上超过匹配换入数量的随机替换；错配问题在已完成的 GQA 上弱于真实问题。TextVQA-OCR 同时下降。覆盖约束限制的是特征代理上的代表关系，不保证答案证据被保住。视觉塔仍然完整执行，全网格投影、成对相似度和剪点上的问题--键乘积都会计入实际成本；各层缓存长度等于该层真正处理过的序列，不能按末层 10 个 token 统一估计。本文没有把上述数量预算写成延迟或显存收益。
